% GENERATED FILE. Edit canonical lessons, then run npm run generate:books.
% canonical-source-hash: fnv1a64:739e26a3186bb988
% canonical-lessons: PA-C363-koshish, PA-C363-usarna, PA-C363-nibhauna, PA-C363-samvarna, PA-C363-tarakki

\chapter{At Work}
\label{ch:pa-a2-at-work}
\hlchaptermodality{363}

\begin{chapteropening}
\textbf{By the end of this chapter:} \emph{I can say what I do and manage at work.}

Complete the last lesson of chapter 363: I can say what I do and manage at work.
\end{chapteropening}

\section[\texorpdfstring{\pa{ਕੋਸ਼ਿਸ਼} (koshish)}{ਕੋਸ਼ਿਸ਼ (koshish)}]{\pa{ਕੋਸ਼ਿਸ਼} (koshish) — an effort; a try}

\label{lesson:PA-C363-koshish}



\begin{quote}
Before the new one: say the Punjabi for to carve, then the Punjabi for to try out.
\end{quote}

\subsection*{You'll want to know: \pa{ਕੋਸ਼ਿਸ਼}}
\textbf{\pa{ਕੋਸ਼ਿਸ਼}} — \emph{koshish} — \textquotedblleft{}an effort; a try\textquotedblright{}.

\begin{quote}
\pa{ਮੈਂ} \pa{ਹਰ} \pa{ਰੋਜ਼} \pa{ਕੋਸ਼ਿਸ਼} \pa{ਕਰਦਾ} \pa{ਹਾਂ।}
\end{quote}

\begin{quote}
\emph{maiṁ har roz koshish kardā hāṁ.}
\end{quote}

\begin{quote}
I make an effort every day.
\end{quote}

A feminine noun.

\begin{cousinweb}
\textbf{\pa{ਕੋਸ਼ਿਸ਼}} (\emph{koshish}) is a loan, usually traced to Persian \emph{kōshish}, \textquotedblleft{}effort\textquotedblright{}.
\end{cousinweb}

\begin{grammarlens}[title={What you've built}]
One more word for this chapter.
\end{grammarlens}

\subsection*{Your turn}
\begin{itemize}
\raggedright
  \item \emph{Say it:} \emph{koshish}
  \item \emph{Say it:} \emph{koshish}, once more
  \item \emph{Say it:} \emph{maiṁ har roz koshish kardā hāṁ.} — I make an effort every day.
  \item \emph{Recall:} say the Punjabi for to carve, then the Punjabi for to try out, then say \emph{koshish} again
\end{itemize}

\begin{tcolorbox}[breakable,colback=teal!4,colframe=teal!35!black,title={Before you move on}]
What does \pa{ਕੋਸ਼ਿਸ਼} mean? (\textquotedblleft{}An effort; a try\textquotedblright{}.) Say it once more.
\end{tcolorbox}
\section[\texorpdfstring{\pa{ਉਸਾਰਨਾ} (usārnā)}{ਉਸਾਰਨਾ (usārnā)}]{\pa{ਉਸਾਰਨਾ} (usārnā) — to put up a building; to build}

\label{lesson:PA-C363-usarna}



\begin{quote}
Before the new one: say the Punjabi for to try out, then the Punjabi for an effort.
\end{quote}

\subsection*{You'll want to know: \pa{ਉਸਾਰਨਾ}}
\textbf{\pa{ਉਸਾਰਨਾ}} — \emph{usārnā} — \textquotedblleft{}to put up a building; to build\textquotedblright{}.

\begin{quote}
\pa{ਮੇਰਾ} \pa{ਪਿਤਾ} \pa{ਪਿੰਡ} \pa{ਵਿੱਚ} \pa{ਘਰ} \pa{ਉਸਾਰਦਾ} \pa{ਹੈ।}
\end{quote}

\begin{quote}
\emph{merā pitā piṇḍ vich ghar usārdā hai.}
\end{quote}

\begin{quote}
My father builds houses in the village.
\end{quote}

Past (I …): \pa{ਮੈਂ} \pa{ਇੱਕ} \pa{ਕਮਰਾ} \pa{ਉਸਾਰਿਆ}, \emph{maiṁ ikk kamrā usāriā}.

\begin{cousinweb}
\textbf{\pa{ਉਸਾਰਨਾ}} (\emph{usārnā}) is the transitive partner of \emph{ussarṇā}, \textquotedblleft{}to rise (of a building)\textquotedblright{}, which goes back to Sanskrit \emph{útsarati}, \textquotedblleft{}hastens away\textquotedblright{}, with a sense \textquotedblleft{}rises\textquotedblright{} that CDIAL infers; CDIAL lists Punjabi \emph{ussarṇā} under it (CDIAL 1874), but not the transitive form itself.
\end{cousinweb}

\begin{grammarlens}[title={What you've built}]
One more word for this chapter.
\end{grammarlens}

\subsection*{Your turn}
\begin{itemize}
\raggedright
  \item \emph{Say it:} \emph{usārnā}
  \item \emph{Say it:} \emph{usārnā}, once more
  \item \emph{Say it:} \emph{merā pitā piṇḍ vich ghar usārdā hai.} — My father builds houses in the village.
  \item \emph{Recall:} say the Punjabi for to try out, then the Punjabi for an effort, then say \emph{usārnā} again
  \item \emph{Write it:} copy \textbf{\pa{ਉਸਾਰਨਾ}} once, checking the whole word
\end{itemize}

\begin{tcolorbox}[breakable,colback=teal!4,colframe=teal!35!black,title={Before you move on}]
What does \pa{ਉਸਾਰਨਾ} mean? (\textquotedblleft{}To put up a building; to build\textquotedblright{}.) Say it once more.
\end{tcolorbox}
\section[\texorpdfstring{\pa{ਨਿਭਾਉਣਾ} (nibhāuṇā)}{ਨਿਭਾਉਣਾ (nibhāuṇā)}]{\pa{ਨਿਭਾਉਣਾ} (nibhāuṇā) — to carry out a duty; to keep a promise}

\label{lesson:PA-C363-nibhauna}



\begin{quote}
Before the new one: say the Punjabi for an effort, then the Punjabi for to put up a building.
\end{quote}

\subsection*{You'll want to know: \pa{ਨਿਭਾਉਣਾ}}
\textbf{\pa{ਨਿਭਾਉਣਾ}} — \emph{nibhāuṇā} — \textquotedblleft{}to carry out a duty; to keep a promise\textquotedblright{}.

\begin{quote}
\pa{ਉਹ} \pa{ਹਮੇਸ਼ਾ} \pa{ਆਪਣਾ} \pa{ਵਾਅਦਾ} \pa{ਨਿਭਾਉਂਦਾ} \pa{ਹੈ।}
\end{quote}

\begin{quote}
\emph{oh hameshā āpṇā vāadā nibhāundā hai.}
\end{quote}

\begin{quote}
He always keeps his promise.
\end{quote}

Past (I …): \pa{ਮੈਂ} \pa{ਆਪਣਾ} \pa{ਵਾਅਦਾ} \pa{ਨਿਭਾਇਆ}, \emph{maiṁ āpṇā vāadā nibhāiā}.

\begin{cousinweb}
\textbf{\pa{ਨਿਭਾਉਣਾ}} (\emph{nibhāuṇā}) is inherited from Sanskrit \emph{nirvahati}, \textquotedblleft{}leads out of, accomplishes\textquotedblright{}; CDIAL lists Punjabi \emph{nibhāuṇā}, \textquotedblleft{}to accomplish\textquotedblright{}, under it (CDIAL 7397).
\end{cousinweb}

\begin{grammarlens}[title={What you've built}]
One more word for this chapter.
\end{grammarlens}

\subsection*{Your turn}
\begin{itemize}
\raggedright
  \item \emph{Say it:} \emph{nibhāuṇā}
  \item \emph{Say it:} \emph{nibhāuṇā}, once more
  \item \emph{Say it:} \emph{oh hameshā āpṇā vāadā nibhāundā hai.} — He always keeps his promise.
  \item \emph{Recall:} say the Punjabi for an effort, then the Punjabi for to put up a building, then say \emph{nibhāuṇā} again
  \item \emph{Write it:} copy \textbf{\pa{ਨਿਭਾਉਣਾ}} once, checking the whole word
\end{itemize}

\begin{tcolorbox}[breakable,colback=teal!4,colframe=teal!35!black,title={Before you move on}]
What does \pa{ਨਿਭਾਉਣਾ} mean? (\textquotedblleft{}To carry out a duty; to keep a promise\textquotedblright{}.) Say it once more.
\end{tcolorbox}
\section[\texorpdfstring{\pa{ਸੰਵਾਰਨਾ} (saṁvārnā)}{ਸੰਵਾਰਨਾ (saṁvārnā)}]{\pa{ਸੰਵਾਰਨਾ} (saṁvārnā) — to tidy up; to set in order}

\label{lesson:PA-C363-samvarna}



\begin{quote}
Before the new one: say the Punjabi for to put up a building, then the Punjabi for to carry out a duty.
\end{quote}

\subsection*{You'll want to know: \pa{ਸੰਵਾਰਨਾ}}
\textbf{\pa{ਸੰਵਾਰਨਾ}} — \emph{saṁvārnā} — \textquotedblleft{}to tidy up; to set in order\textquotedblright{}.

\begin{quote}
\pa{ਮੈਂ} \pa{ਸਵੇਰੇ} \pa{ਦੁਕਾਨ} \pa{ਸੰਵਾਰਦਾ} \pa{ਹਾਂ।}
\end{quote}

\begin{quote}
\emph{maiṁ savere dukān saṁvārdā hāṁ.}
\end{quote}

\begin{quote}
I set the shop in order in the morning.
\end{quote}

Past (I …): \pa{ਮੈਂ} \pa{ਕਮਰਾ} \pa{ਸੰਵਾਰਿਆ}, \emph{maiṁ kamrā saṁvāriā}.

\begin{cousinweb}
\textbf{\pa{ਸੰਵਾਰਨਾ}} (\emph{saṁvārnā}) is inherited from a reconstructed \emph{*śodhakāra}, \textquotedblleft{}making clean\textquotedblright{}; CDIAL lists Punjabi \emph{savārṇā}, \textquotedblleft{}to adorn, improve, prepare, adjust\textquotedblright{}, under it (CDIAL 12628); the spelling with the nasal is not listed there.
\end{cousinweb}

\begin{grammarlens}[title={What you've built}]
One more word for this chapter.
\end{grammarlens}

\subsection*{Your turn}
\begin{itemize}
\raggedright
  \item \emph{Say it:} \emph{saṁvārnā}
  \item \emph{Say it:} \emph{saṁvārnā}, once more
  \item \emph{Say it:} \emph{maiṁ savere dukān saṁvārdā hāṁ.} — I set the shop in order in the morning.
  \item \emph{Recall:} say the Punjabi for to put up a building, then the Punjabi for to carry out a duty, then say \emph{saṁvārnā} again
  \item \emph{Write it:} copy \textbf{\pa{ਸੰਵਾਰਨਾ}} once, checking the whole word
\end{itemize}

\begin{tcolorbox}[breakable,colback=teal!4,colframe=teal!35!black,title={Before you move on}]
What does \pa{ਸੰਵਾਰਨਾ} mean? (\textquotedblleft{}To tidy up; to set in order\textquotedblright{}.) Say it once more.
\end{tcolorbox}
\section[\texorpdfstring{\pa{ਤਰੱਕੀ} (tarakkī)}{ਤਰੱਕੀ (tarakkī)}]{\pa{ਤਰੱਕੀ} (tarakkī) — a promotion; progress}

\label{lesson:PA-C363-tarakki}



\begin{quote}
Before the new one: say the Punjabi for to carry out a duty, then the Punjabi for to set in order.
\end{quote}

\subsection*{You'll want to know: \pa{ਤਰੱਕੀ}}
\textbf{\pa{ਤਰੱਕੀ}} — \emph{tarakkī} — \textquotedblleft{}a promotion; progress\textquotedblright{}.

\begin{quote}
\pa{ਮੇਰੇ} \pa{ਪਿਤਾ} \pa{ਨੂੰ} \pa{ਦਫ਼ਤਰ} \pa{ਵਿੱਚ} \pa{ਤਰੱਕੀ} \pa{ਮਿਲੀ।}
\end{quote}

\begin{quote}
\emph{mere pitā nū̃ daftar vich tarakkī milī.}
\end{quote}

\begin{quote}
My father got a promotion at the office.
\end{quote}

A feminine noun.

\begin{cousinweb}
\textbf{\pa{ਤਰੱਕੀ}} (\emph{tarakkī}) is a loan, usually traced through Persian to Arabic \emph{taraqqī}, \textquotedblleft{}rising\textquotedblright{}.
\end{cousinweb}

\begin{grammarlens}[title={What you've built}]
One more word for this chapter.
\end{grammarlens}

\subsection*{Your turn}
\begin{itemize}
\raggedright
  \item \emph{Say it:} \emph{tarakkī}
  \item \emph{Say it:} \emph{tarakkī}, once more
  \item \emph{Say it:} \emph{mere pitā nū̃ daftar vich tarakkī milī.} — My father got a promotion at the office.
  \item \emph{Recall:} say the Punjabi for to carry out a duty, then the Punjabi for to set in order, then say \emph{tarakkī} again
  \item \emph{Write it:} copy \textbf{\pa{ਤਰੱਕੀ}} once, checking the whole word
\end{itemize}

\begin{tcolorbox}[breakable,colback=teal!4,colframe=teal!35!black,title={Before you move on}]
What does \pa{ਤਰੱਕੀ} mean? (\textquotedblleft{}A promotion; progress\textquotedblright{}.) Say it once more.
\end{tcolorbox}
